\documentclass[11pt]{article}
\usepackage[a4paper,margin=27mm]{geometry}
\usepackage{amsmath,amssymb,hyperref}
\title{The cyclic coefficient test inside the author's function class}
\author{}\date{20 September 2026}
\begin{document}\maketitle

This finite check supplements SA10--SA11 of
\href{https://github.com/KokunoYumeto/zeta-function-research-reader/blob/5992ad50c0f2a06921f1fcaded7b4e38097c0739/workbenches/splitzero-tandem/continuations/20260920-original-kernel-mixed-determinants/providers/NATIVE_SPECTRAL_ACTION.tex\#L1}{\texttt{NATIVE\_SPECTRAL\_ACTION.tex}}. Its purpose is to test
the coefficient in the statement of Connes--van Suijlekom,
\emph{Quadratic Forms, Real Zeros and Echoes of the Spectral Action},
\href{https://arxiv.org/abs/2511.23257v1}{arXiv:2511.23257v1},
Section \texttt{sect:sa}, Lemma \texttt{lem:sa}, with a
function actually belonging to the author's stated class.
It does not substitute an unbounded polynomial for that class.

Use precisely the Fourier convention
\[
f(x)=\int_{\mathbb R}\widehat f(\xi)e^{i\xi x}\,d\xi,
\qquad
f(x)=x^2e^{-x^2},\qquad
\widehat f(\xi)=\frac1{2\sqrt\pi}
\left(\frac12-\frac{\xi^2}{4}\right)e^{-\xi^2/4}.
\tag{FC1}
\]
Here is a direct verification of the constants.
The Gaussian integral
$J(\xi)=\int_{\mathbb R}e^{-x^2-i\xi x}\,dx$ satisfies
$J'(\xi)=-\xi J(\xi)/2$ by integration by parts, with
$J(0)=\sqrt\pi$ from squaring the real Gaussian integral
and using polar coordinates. Therefore
$J(\xi)=\sqrt\pi e^{-\xi^2/4}$.
Differentiating twice in $\xi$ under the absolutely
convergent integral shows that the Fourier transform
of $x^2e^{-x^2}$ is $-J''(\xi)$.
The factor $1/(2\pi)$ in the inverse convention gives
exactly (FC1). Equivalently insert the displayed
Gaussian integral into (FC1) and differentiate with
respect to $x$; this verifies its inverse directly.

For every integer $n\geq0$, integration by parts gives
$\widehat{f^{(n)}}(\xi)=(i\xi)^n\widehat f(\xi)$.
Set
\[
K=\int_{\mathbb R}e^{|\xi|}|\widehat f(\xi)|\,d\xi<\infty,
\qquad C=\max(1,K).
\tag{FC2}
\]
The finiteness follows directly from (FC1), since a
quadratic polynomial times $e^{|\xi|-\xi^2/4}$ is
integrable. The exponential series gives
$|\xi|^n\leq n!e^{|\xi|}$. Consequently
\[
\|\widehat{f^{(n)}}\|_1
\leq K n!\leq C^{n+1}n!.
\tag{FC3}
\]
Thus $f$ belongs to the source's class $\mathcal E^0$:
the auxiliary powers $(x-i)^m$ required there have
only $m=0$ when its summability index is $s=0$.

Now take the one-dimensional Hilbert space, $D=0$ and
$R=1$. Its identity $(D-i)^0$ is trace class, $D$ is
selfadjoint and $R$ is bounded, so these are exactly
data admitted by that lemma with $s=0$. Direct evaluation
gives
\[
\operatorname{tr}f(D+tR)=t^2e^{-t^2},\qquad
\left.\frac{d^2}{dt^2}\operatorname{tr}f(D+tR)\right|_{t=0}
=2.
\tag{FC4}
\]
The sole cyclic product of matrix entries is $1$ and
the repeated divided difference is
$f'[0,0]=f''(0)=2$.
The displayed $n!$ coefficient in the lemma's statement
would give $2!\,f'[0,0]=4$. The coefficient
$(n-1)!$ proved in SA10 gives the actual value $2$.
The source proof's repeated-node expression instead
gives $2!\,f[0,0,0]=2!\,f''(0)/2!=2$, also correct.

This establishes the coefficient correction within the
stated source class itself. It does not change the
source's subsequent second-derivative formula or the
finite native identities SA1--SA43.
The original author LaTeX read for this check is
\texttt{Araki-final-oct25.tex}, Section \texttt{sect:sa},
SHA-256
\texttt{0357a6123b7b3a2ba2e058458f5267eef1c4ad6abb051af559015b86ae2aa111}.
No PDF or other edition was used.
\end{document}
